\documentclass[11pt,a4paper]{article}

\usepackage[utf8]{inputenc}
\usepackage[T2A,T1]{fontenc}
\usepackage[english,russian]{babel}
\usepackage[margin=1.6cm]{geometry}
\usepackage{microtype}
\usepackage{booktabs}
\usepackage{amsmath,amssymb}
\usepackage{xcolor}
\usepackage{hyperref}
\usepackage{enumitem}
\usepackage{titlesec}
\usepackage{titling}

\hypersetup{colorlinks=true, linkcolor=blue!50!black, urlcolor=blue!50!black}

\titleformat{\section}{\normalfont\large\bfseries}{\thesection}{0.6em}{}
\titleformat{\subsection}{\normalfont\normalsize\bfseries}{\thesubsection}{0.5em}{}
\titlespacing*{\section}{0pt}{4pt}{0pt}
\titlespacing*{\subsection}{0pt}{3pt}{0pt}

\setlist[itemize]{leftmargin=1.4em,topsep=2pt,itemsep=1pt,parsep=1pt}
\setlist[enumerate]{leftmargin=1.6em,topsep=2pt,itemsep=2pt,parsep=1pt}
\setlength{\parskip}{1pt}
\setlength{\parindent}{0pt}

\setlength{\droptitle}{-2.5em}
\pretitle{\begin{flushleft}\Large\bfseries}
\posttitle{\end{flushleft}\vspace{-0.6em}}
\preauthor{}\postauthor{}
\predate{}\postdate{}

\title{SPC: self-play-критик для оценки рассуждений. Обзор статьи и эксперимент}
\author{}
\date{}

\pagestyle{empty}

\begin{document}
\maketitle

\noindent\small Статья: Chen et al., \textit{SPC: Evolving Self-Play Critic via Adversarial Games for LLM Reasoning}, NeurIPS 2025, \href{https://arxiv.org/abs/2504.19162}{arXiv:2504.19162}.\normalsize

\section*{Оригинальная статья}

Авторы ставят перед собой цель — научить модель находить ошибки в пошаговых математических решениях без ручной разметки каждого шага. Они берут две копии Qwen2.5-7B-Instruct, дообучают их на размеченных данных (стартовое дообучение, SFT) и сталкивают в игре двух ролей.

\begin{itemize}
  \item \textbf{Генератор-обманщик} переписывает один шаг корректного решения в неверный, так чтобы тот не выдавал себя, и пытается обмануть критика.
  \item \textbf{Критик}, видя задачу, предыдущие шаги и проверяемый шаг, после короткого разбора отвечает \texttt{Correct} или \texttt{Incorrect}.
\end{itemize}

За раунд разыгрываются десятки тысяч партий. Награду определяет исход — очко достаётся тому, кто переиграл, — и по ней обе модели дообучаются через PPO: генератор учится прятать ошибки тоньше, критик находить их точнее. После SFT разметка больше не нужна, работает self-play.

Авторы выложили две версии критика — Critic-0 (после SFT) и Critic-2 (после двух раундов); генератор не опубликован.

\subsection*{Задачи}

Центральная задача — построить надёжного пошагового критика без ручной разметки шагов. Пошаговая оценка нужна, чтобы награждать модель за каждый шаг, а не только за финальный ответ. Сильные критики такого рода (process-reward-модели, PRM) обучаются на тысячах размеченных людьми шагов (PRM800K от OpenAI), и эта разметка не масштабируется. Отсюда две цели: порождать обучающий сигнал автоматически и не дать ему упрощаться по мере роста критика. Целью был критик, который через self-play догоняет на ProcessBench PRM.

\subsection*{Сильные стороны}

Разметка нужна только для SFT — дальше процесс обходится без человека. Награда за каждый шаг плотнее награды за финальный ответ, поэтому обучение сходится быстрее. Архитектура простая: одна 7B-модель в обеих ролях, без отдельной reward model и value-функции. При этом результат подтвердил эффективность — на ProcessBench Critic-2 даёт около $78\%$ F1, на уровне PRM, обученных на размеченных данных.

\subsection*{Слабые стороны}

Критик ловит лишь те ошибки, которые умеет вставлять текущий генератор, — это искусственное распределение, тогда как ProcessBench собран из настоящих ошибок моделей. Расхождение train/test заложено в самом методе, но в статье не обсуждается.

В релизе две версии критика из трёх, а генератор закрыт — нельзя ни воспроизвести его сторону обучения, ни заглянуть за второй раунд.

Вся оценка сведена к одному числу — ProcessBench F1, по которому не видно, что изменила самоигра: вырос recall или просто упали ложные срабатывания на корректных шагах. Это я и проверил.

\section*{Гипотеза эксперимента}

В каждом обучающем примере критик видит шаг с подложенной ошибкой, так что после многих раундов его prior должен смещаться к ответу <<Incorrect>>. Предсказание: Critic-2 точнее ловит настоящие ошибки (выше recall), но чаще принимает корректные шаги за неверные (ниже accuracy). Поэтому я раскладываю общую метрику на recall на ошибочных шагах и accuracy на корректных.

\section*{Эксперимент}

Беру обе открытые версии (\texttt{judge/SPC-Critic-0}, \texttt{judge/SPC-Critic-2}) и первые 40 объектов каждого подмножества ProcessBench (\texttt{gsm8k}, \texttt{math}, \texttt{olympiadbench}, \texttt{omnimath}). Каждый разворачиваю в пошаговые пробы по схеме SPC — 296 проб на критика: 136 корректных шагов, 160 с ошибкой. Оба критика идут по одному набору, поэтому сравнение парное. Промпт и формат ответа — из репозитория авторов, декодирование жадное ($T{=}0$), \texttt{max\_new\_tokens=512}, инференс через vLLM. Отступления от оригинала: ProcessBench пересобран из сырого HF-набора, выборка $N{=}296$ против ${\sim}3400$.

\begin{center}
\begin{tabular}{lrrr}
\toprule
метрика & Critic-0 (SFT) & Critic-2 (RL) & $\Delta$ \\
\midrule
recall на ошибочных шагах             & 0.731 & 0.711          & $-0.020$ \\
специфичность на корректных шагах     & 0.674 & 0.859          & $+0.185$ \\
ProcessBench F1 (гарм. среднее, headline) & 0.701 & \textbf{0.778} & $+0.077$ \\
арифм. среднее (balanced accuracy)    & 0.705 & 0.781          & $+0.076$ \\
\bottomrule
\end{tabular}
\end{center}

ProcessBench F1 у Critic-2 — $0.778$ против $0.777$ в статье: расхождение меньше $0.1$ пункта на выборке в 12 раз меньшей. Пайплайн воспроизводит авторские цифры.

\subsection*{Результаты}

По гипотезе recall должен был вырасти, а специфичность — упасть. Вышло наоборот: recall чуть снизился (на 2 пункта), а специфичность (точность на корректных шагах) выросла на 18.5. Ниже сравнение предсказаний C0 и C2:

\begin{center}
\begin{tabular}{lrrr}
\toprule
истинный класс шага & C0 $<$ C2 & C0 $>$ C2 & $\sum$ \\
\midrule
шаги без ошибки ($n{=}136$) & 31 & 6  & $+25$ \\
шаги с ошибкой ($n{=}160$)  & 22 & 26 & $-4$  \\
\bottomrule
\end{tabular}
\end{center}

\smallskip
{\small C0 — Critic-0, C2 — Critic-2; «C0 $<$ C2» — C0 ошибся, а C2 ответил верно; «C0 $>$ C2» — наоборот; $\sum$ $=$ (стало лучше) $-$ (стало хуже).}

\smallskip
На корректных шагах второй раунд даёт сильный прирост качества ($+25$), на ошибочных — немного ломает ($-4$). Значит, прирост даёт не лучшая детекция ошибок, а то, что Critic-2 перестал ошибочно отвергать корректные шаги: критик стал \textit{снисходительнее}, и гипотеза не подтвердилась.

\subsection*{Ограничения}

\begin{enumerate}
  \item 296 проб против ${\sim}3400$ в полном ProcessBench; доверительных интервалов нет, поэтому небольшие сдвиги (особенно recall, $-2$ пункта) могут лежать в пределах шума.
  \item ProcessBench развёрнут в пошаговые пробы моей схемой, а не оригинальными eval-данными авторов; абсолютные числа могут смещаться — надёжно лишь сравнение Critic-0 vs Critic-2.
  \item В релизе лишь Critic-0 и Critic-2, промежуточные раунды недоступны, поэтому вывод о направлении тренда опирается всего на две точки.
  \item Жадное декодирование без повторов и единственный домен (математика ProcessBench); разброс по сэмплированию и перенос на другие задачи не проверялись.
\end{enumerate}

\section*{Следующие шаги}

Эксперимент воспроизвёл число из статьи, но вскрыл его природу: прирост даёт рост специфичности (критик стал мягче), а не более чуткая детекция ошибок. Пока это в плюс, но направление опасное — с ростом числа итераций критик может начать терять настоящие ошибки быстрее, чем выигрывает на корректных. Отсюда главный вопрос: монотонен ли этот тренд, есть ли раунд, где самоигра переусердствует, и почему критик смягчается.

Ответить на него можно следующими экспериментами:

\begin{enumerate}
  \item Прогнать промежуточные раунды (Critic-1, Critic-3, \dots) — их придётся дообучить, поскольку в релизе только Critic-0 и Critic-2: подтвердится ли тренд и есть ли точка, где recall начинает падать.
  \item Разложить пропущенные ошибки по типам (описка, неверная формула, путаница со знаком, пропущенный шаг) — это покажет, куда сместилось распределение генератора.
\end{enumerate}

\section*{Источники}

{\footnotesize
\begin{enumerate}[label={[\arabic*]},leftmargin=2.2em,topsep=2pt,itemsep=1pt]
  \item J.~Chen и др. \textit{SPC: Evolving Self-Play Critic via Adversarial Games for LLM Reasoning}. NeurIPS 2025. arXiv:2504.19162.
  \item Qwen Team. \textit{Qwen2.5 Technical Report}. 2024. arXiv:2412.15115.
  \item J.~Schulman и др. \textit{Proximal Policy Optimization Algorithms}. 2017. arXiv:1707.06347.
  \item H.~Lightman и др. \textit{Let's Verify Step by Step}. 2023. arXiv:2305.20050. (датасет PRM800K)
  \item C.~Zheng и др. \textit{ProcessBench: Identifying Process Errors in Mathematical Reasoning}. 2024. arXiv:2412.06559.
  \item W.~Kwon и др. \textit{Efficient Memory Management for LLM Serving with PagedAttention} (vLLM). SOSP 2023. arXiv:2309.06180.
\end{enumerate}
}

\end{document}